In this work, we have described a general \emph{gauging measurement} protocol by which Clifford logicals of generic QLDPC codes may be measured fault-tolerantly. First, the Clifford operator is gauged, deforming the initial code into a non-Abelian Clifford stabilizer code. Next, the system is deformed back onto the initial code and the input code state is projected onto an eigenstate of the target Clifford logical. We have proven that the full measurement protocol is fault tolerant, given an input qubit Pauli stabilizer code with an order-$2$ transversal Clifford. This treatment of fault-tolerance for non-Abelian QLDPC codes goes beyond the heuristic discussions that have appeared previously in the literature. 

The logical Clifford magic states prepared by our protocol applied to general CSS QLDPC codes may be used to perform non-Clifford operations, either by conversion to a conventional $\ket{\mathsf{Tof}}$ or $\ket{CCZ}$ state, or through a gadget such as those we have introduced in Sec.~\ref{sec:applications}. In the applications of the gauging measurement protocol we have discussed, we were able to reduce the magic state cost for several useful algorithmic subroutines compared to prior methods. We have also discussed several examples of codes resulting from gauging transversal gates in Sec~\ref{sec:examples}.

Our work presents an approach to fault-tolerant magic state preparation, that is distinct from distillation, cultivation, and transversal non-Clifford strategies. Our method only requires the existence of a transversal Clifford operator, and as such the gauging measurement enables us to prepare magic states on a wide array of codes, including generic CSS codes. We avoid the need for costly post-selection by leveraging the distance and LDPC nature of the intermediate gauged code. 

The appearance of Clifford codes in our protocol introduces new complications as compared to Pauli measurement schemes. The stabilizer group of a Clifford stabilizer code may not commute outside of the codespace, detectors become randomized in the presence of certain faults. These randomized detectors no longer provide useful information about the fault history. We have developed new theoretical tools to analyze these complications; in particular, we have developed a framework through which the fault-tolerance properties of Clifford LDPC codes may be analyzed in the presence of such randomized detectors.

Our work on the fault-tolerance of procedures for logical Clifford measurement on QLDPC codes has uncovered a number of new directions which we discuss here, and leave to future exploration. 

\textbf{New useful magic states? }We have presented some applications of Clifford gauging measurements aside from magic state preparation; exploring the landscape of similar ``gauge-able'' gadgets is a promising direction. In particular, it would be interesting to better understand the kinds of structures in circuits that can be made more efficient through a gauging operation. The $\SWAP$ network example we discussed in this work indicates that algorithms containing many gates conditioned on one ancilla can be compressed in a straightforward manner. Are there other common circuit motifs one can produce useful resource states for via gauging?

\textbf{Going beyond qubits? }Currently, our fault distance proof is limited to the case of gauging $\mathbb{Z}_2$ Clifford gates possessed by qubit stabilizer codes. While this encompasses most relevant magic states and Clifford gates, there are interesting Clifford gates that do not fit these requirements. As an explicit example, take the order-3 facet gate $HS$ on the self-dual BB code we considered in Sec.~\ref{sec:examples}. Or we may go beyond the Clifford hierarchy and consider transversal gates  applied to \emph{qudit} codes, see the example of gauging charge conjugation symmetry on a qutrit BB code in Sec.~\ref{sec:examples}. In both cases, the resulting algebra of gauged code stabilizers has a more complex structure than the $\mathbb{Z}_2$ case. Specifically, there are some non-Abelian faults that drop their own detectors, making the spacetime cleaning procedure ambiguous. However, there may be hope to construct a decoder for such codes explicitly by modifying techniques leveraged for similar topological codes~\cite{lyons2026}. 


\textbf{Exotic phases of matter? }One may ask, what is the general structure of the non-Pauli codes produced by gauging Clifford logical operators? Can they be organized into well-defined phases of matter, potentially through a topological defect network picture? Can we produce more exotic models by gauging nonabelian group symmetries, or even Hopf algebra symmetries? Can these techniques and constructions be extended to continuous variable systems?


\textbf{Decoding? }We have proved that many Clifford gauging protocols for stabilizer qLDPC codes have nontrivial spacetime distance. Putting these protocols into practice will require some sort of decoding while we remain in the intermediate Clifford stabilizer code. We leave a more detailed treatment of decoding strategies for the Clifford LDPC codes introduced here to future work. However, there are a few directions we believe seem promising to explore. In particular, decoders known to work well for expander codes, like small-set flip (SSF) \cite{Leverrier_2015,Fawzi_2018_2,Fawzi_2018} and windowed BP-OSD \cite{Panteleev_2021, roffe2020, huang2023improvednoisysyndromedecoding, gong2024lowlatencyiterativedecodingqldpc} could be generalized in a straightforward way to $\mathbb{Z}_2$ codes with $\mathbb{Z}_2$ gauged Clifford symmetries. We note that just-in-time decoders have been proven to work for non-Abelian \emph{topological} codes \cite{davydova2025universal, lyons2026}. These decoders possess a renormalization group-type structure which is tailored to the finite-dimensional nature of topological codes, but it remains to be seen whether they are suitable for expander geometry. However, one could envision combining the just-in-time structure with a better heuristic decoder like BP-OSD. We predict that it will be necessary to incorporate a window-decoding strategy alongside the BP-OSD, since non-Abelian faults that persist for any extensive period of time can hide an extensive number of Abelian faults, and thus reduce the spacetime distance of the procedure. 

\textbf{Numerical simulation? }After constructing a viable decoder, it remains to be seen what the overall resource overhead for the gauging measurement protocol looks like in practice. Recent work \cite{delafuente2026highthresholddecodingnonpaulicodes} has demonstrated that just-in-time decoding for a topological Clifford code supporting $D(D_4)$ anyons can perform comparably to the surface code under circuit-level noise. If we conjecture that the threshold for the Clifford measurement protocol is similar to the memory threshold of the underlying Clifford code, this suggests that the performance of gauging-based methods may be comparable to stabilizer ones. Recently introduced efficient Clifford code simulators may be leveraged to make the necessary numerics possible ~\cite{takada2026exactefficientsimulationnoisy, Surti_2026,delafuente2026efficientlyestimatingfailurerates}. 

\textbf{Fast, High-rate surgery? }We have not attempted to optimize any overhead or resource-costs for the protocol introduced here. It is natural to consider how to implement fast Clifford surgery in the presence of many logical representatives of the Clifford transversal in question, see Ref.~\cite{williamson2026fast}. Determining the properties that the input code and the gauging graph need to satisfy to allow for fault-tolerant fast surgery, as well as identifying promising code candidates, is a topic for future work. It is also interesting to consider how we can improve the rate of magic state preparation through parallel gauging methods: we presented a protocol to generate a set of $k$ disjoint logical magic states on a high-rate code with an unaddressable transversal gate in an $O(k)$ number of gauging measurements, performed serially. Given a code with addressable Clifford gates, it may be possible to create $O(k)$ magic states in parallel following a procedure similar to Ref.~\cite{cowtan2025parallel,Zhang_2025}, reducing time overhead. Incorporating Clifford gauging measurement into existing architectures, such as the fixed-connectivity extractor architecture \cite{he2025extractorsqldpcarchitecturesefficient} is also interesting. In this setting, we imagine that the base code has a transversal Clifford operator, such as a $\CNOT$ between codeblocks, that one can gauge to prepare magic states or perform a particular magic gate. All extra gauge qubits and qubit connectivity required to measure $\CNOT$ can be fixed in advance for each disjoint pair of codeblocks. Consequently, a pair of codeblocks would be able to dynamically change their role in the compilation from computational logical processing units to magic state factories. More generally, extractors can be generalized to support Clifford measurements on a code block by constructing an ancilla based on the graph desiderata presented above for the whole code (or a targeted subregion) rather than the Pauli measurement graph desiderata.




